\chapter{Fifth Disease (Erythema Infectiosum)}
\index{Fifth Disease (Erythema Infectiosum)}
\label{sec:Fifth Disease (Erythema Infectiosum)}

\section{Overview}
Fifth disease, also known as erythema infectiosum, is a common childhood viral exanthem caused by parvovirus B19. It is characterized by a distinctive slapped cheek rash followed by a lacy reticular rash on the trunk and extremities. The infection is typically mild and self-limited in children but can cause significant complications in adults, immunocompromised individuals, and during pregnancy.
\section{Classification}

\begin{description}[font=\normalfont\bfseries,leftmargin=3em,labelwidth=2.5em]
  \item[Cause] Infectious (Viral)
  \item[Organ System] Pediatric / Infectious
  \item[Pathophysiology] Parvovirus B19 infection causing lytic infection of erythroid progenitor cells
  \item[Duration] Acute (1--3 weeks)
  \item[Onset] Sudden
  \item[Mode of Transmission] Respiratory droplets; bloodborne; vertical (mother to fetus)
  \item[Clinical Course] Acute, self-limited
  \item[Severity] Mild
  \item[Extent of Disease] Systemic
  \item[Age Group] Children (peak 4--10 years)
  \item[Epidemiology] Very common; 40--60\% of adults seropositive; epidemic cycles every 3--6 years
  \item[ICD-11 Code] 1D60.0
\end{description}

\section{Causes}
Parvovirus B19, a small single-stranded DNA virus, infects erythroid progenitor cells via the P blood group antigen (globoside) receptor, leading to temporary arrest of red blood cell production. The rash and joint symptoms are immune-mediated, appearing after the development of IgM antibodies. Transmission occurs through respiratory secretions, blood products, and vertically from mother to fetus.

\section{Risk Factors}

\begin{itemize}[noitemsep]
  \item Age 4--10 years (peak incidence)
  \item Contact with infected children (parents, teachers, healthcare workers)
  \item Outdoor activities and school attendance
  \item Immunosuppression (chronic hemolytic anemia, HIV, transplant)
  \item Pregnancy (risk of fetal hydrops)
  \item Lack of prior exposure (seronegative individuals)
\end{itemize}

\section{Signs and Symptoms}

\subsection{Early symptoms}
\begin{itemize}[noitemsep]
  \item Early manifestations of fifth disease (erythema infectiosum) may be subtle and nonspecific
\end{itemize}

\subsection{Common symptoms}
\begin{itemize}[noitemsep]
  \item \subsection{Early symptoms}
  \item \subsection{Common symptoms}
  \item \textbf{Common symptoms:}
  \item Prodrome: low-grade fever, headache, malaise, myalgia (mild or absent in children)
  \item Slapped cheek rash: bright red, confluent erythema on both cheeks, sparing the nose and perioral area
  \item Reticular rash: lacy, erythematous maculopapular rash on trunk and extremities
  \item Rash may recur with heat, exercise, or sun exposure for several weeks
  \item Arthralgia and arthritis (more common in adults, especially women)
  \item Pain or discomfort in the affected area
  \item Difficulty performing daily activities
\end{itemize}

\subsection{Less common symptoms}
\begin{itemize}[noitemsep]
  \item Atypical or less common presentations of fifth disease (erythema infectiosum) may occur
\end{itemize}

\subsection{Emergency warning signs}
\begin{itemize}[noitemsep]
  \item Severe or worsening symptoms of fifth disease (erythema infectiosum) require immediate medical evaluation
\end{itemize}
\section{Progression and Stages}
After an incubation period of 4--14 days, the prodromal phase (if present) lasts 2--3 days. The slapped cheek rash appears first, followed 1--4 days later by the reticular rash that spreads to the trunk and limbs. The rash typically lasts 1--3 weeks but may recur with triggers. Arthralgia can persist for weeks to months, particularly in adults.

\section{Complications}

\begin{itemize}[noitemsep]
  \item Aplastic crisis in patients with hemolytic anemias (sickle cell disease, hereditary spherocytosis)
  \item Chronic pure red cell aplasia in immunocompromised patients
  \item Hydrops fetalis and fetal loss in pregnancy (infection in first 20 weeks)
  \item Chronic arthropathy in adults
  \item Myocarditis and hepatitis (rare)
  \item Reduced quality of life
  \item Emotional and psychological effects
\end{itemize}

\section{Diagnosis}

\begin{itemize}[noitemsep]
  \item Clinical diagnosis based on characteristic slapped cheek rash
  \item Parvovirus B19 IgM serology confirms recent infection
  \item PCR for parvovirus B19 DNA in immunocompromised or fetal cases
  \item Complete blood count showing reticulocytopenia during aplastic crisis
  \item Prenatal ultrasound for fetal hydrops evaluation
\end{itemize}

\section{Prevention}

\begin{itemize}[noitemsep]
  \item Good hand hygiene and respiratory etiquette
  \item Exclusion from school during prodromal phase (before rash appears; no longer infectious once rash develops)
  \item Pregnant healthcare workers should avoid contact with known cases
  \item No vaccine available currently
  \item Go for regular check-ups so any problems are caught early
\end{itemize}

\section{Treatment Options}

\begin{itemize}[noitemsep]
  \item Supportive care for uncomplicated cases
  \item Antipyretics for fever and joint pain (acetaminophen, ibuprofen, NSAIDs)
  \item Blood transfusion for aplastic crisis
  \item Intravenous immunoglobulin (IVIG) for chronic infection in immunocompromised patients
  \item Intrauterine transfusion for severe fetal anemia and hydrops
  \item Lifestyle changes and self-care
\end{itemize}

\section{Living with It}

\begin{itemize}[noitemsep]
  \item Follow your treatment plan as prescribed
  \item Keep up with regular appointments with your healthcare provider
  \item Adjust your daily habits as needed
  \item Reach out to family, friends, or support groups for help
  \item Keep learning about your condition and how to manage it
\end{itemize}

\section{Outlook}
The prognosis for fifth disease in healthy children is excellent, with full recovery without sequelae. Adults may experience prolonged arthralgia but generally recover fully. The most serious complications occur in patients with underlying hemolytic anemias (aplastic crisis), immunocompromised states, and during pregnancy. With appropriate management, outcomes are generally good even in these higher-risk groups.
